%! TeX program = lualatex
\documentclass[tikz, crop]{standalone}
\usepackage{amsmath, amssymb}

\input{../typography.tex.preamble}
\input{../tikz.tex.preamble}


\pgfplotsset{
  implicit curve/.style={
    basic curve,
    enlargelimits={abs=0.5},
    xmin=-5, xmax=5, ymin=-5, ymax=5,
    width={200pt},
    xtick distance={2},
    ytick distance={2},
    xlabel={\empty}, ylabel={\empty},
    samples=500,
  }
}

\begin{document}
\begin{tikzpicture}
  \begin{axis}[implicit curve]
    \addplot[very thick, domain=-5:5] {sqrt(25 - x^2)};
    \addplot[very thick, domain=-5:5] {-sqrt(25 - x^2)};
  \end{axis}
\end{tikzpicture}

\begin{tikzpicture}
  \begin{axis}[implicit curve]
    \addplot[very thick, domain=-5:5] {sqrt(25 - x^2)};
    \draw[teal, fill] (axis cs:3,4) circle[radius=0.2] node[above right, teal] {\((3,4)\)};
  \end{axis}
\end{tikzpicture}


\begin{tikzpicture}
  \begin{axis}[implicit curve]
    \addplot[very thick, domain=-5:5] {-sqrt(25 - x^2)};
    % \draw[teal, fill] (axis cs:-4,-3) circle[radius=0.2] node[right, teal] {\((-4,-3)\)};
  \end{axis}
\end{tikzpicture}

\begin{tikzpicture}
  \begin{axis}[implicit curve]
    \addplot[very thick, domain=-4:-5] {sqrt(25 - x^2)};
    \addplot[very thick, domain=-4:-5] {-sqrt(25 - x^2)};
    \addplot[very thick, domain=4:5] {sqrt(25 - x^2)};
    \addplot[very thick, domain=4:5] {-sqrt(25 - x^2)};
    \draw[magenta, fill] (axis cs:5,0) circle[radius=0.2] node[above left, supp] {\((5,0)\)};
    \draw[magenta, fill] (axis cs:-5,0) circle[radius=0.2] node[above right, supp] {\((-5,0)\)};
  \end{axis}
\end{tikzpicture}

\begin{tikzpicture}
  \begin{axis}[xtick={\empty}, ytick={\empty}, ymin=-15, ymax=15, width=3in, title={Graph for \(f\)}]
    \addplot[main, very thick, domain=-2:2, axis equal image=false, samples=1000, smooth] {(x^3 + 2*x + 3)};
    % \fill (axis cs:1,6/15*2) circle[radius=0.05];
    \addplot[dashed] coordinates { (1,6)  (1,0) };
    \addplot[dashed] coordinates { (0,6)  (1,6)};
  \end{axis}
\end{tikzpicture}

\begin{tikzpicture}
  \begin{axis}[xtick={\empty}, ytick={\empty}, xmin=-15, xmax=15, width=3in, title={Graph for \(f^{-1}\)}]
    \addplot[supp, very thick, domain=-9:15, axis equal image=false, samples=1000, smooth] {(2*(2/3)^(1/3))/(27 - 9*x + sqrt(3)*sqrt(275 - 162*x + 27*x^2))^(1/3) - (27 - 9*x + sqrt(3)*sqrt(275 - 162*x + 27*x^2))^(1/3)/(2^(1/3)*3^(2/3))};
    % \fill (axis cs:1,6/15*2) circle[radius=0.05];
    \addplot[dashed] coordinates { (6,1)  (0,1) };
    \addplot[dashed] coordinates { (6,0)  (6,1)};
  \end{axis}
\end{tikzpicture}

% Pages 7 and 8: x^2 - xy + y^2 = 1, whole and zoomed in near
% (-1/sqrt(3), 1/sqrt(3)).  The curve is the ellipse traced by
% (cos t - sin t/sqrt(3), cos t + sin t/sqrt(3)); the point is t = pi/2.  Both
% windows are square and drawn at the same size, so the zoom is a true 5x
% magnification.
\pgfmathsetmacro{\zoomx}{-1/sqrt(3)}
\pgfmathsetmacro{\zoomy}{1/sqrt(3)}
\pgfmathsetmacro{\zoomr}{0.3}  % half the side of the zoom window
\pgfmathsetmacro{\zoomxmin}{\zoomx - \zoomr}
\pgfmathsetmacro{\zoomxmax}{\zoomx + \zoomr}
\pgfmathsetmacro{\zoomymin}{\zoomy - \zoomr}
\pgfmathsetmacro{\zoomymax}{\zoomy + \zoomr}

% Q = (-2/sqrt(3), -1/sqrt(3)), t = 5pi/6: the leftmost point, where the
% tangent line is vertical.  Zoomed in on page 9, and boxed in supp on page 7.
\pgfmathsetmacro{\vertx}{-2/sqrt(3)}
\pgfmathsetmacro{\verty}{-1/sqrt(3)}
\pgfmathsetmacro{\vertxmin}{\vertx - \zoomr}
\pgfmathsetmacro{\vertxmax}{\vertx + \zoomr}
\pgfmathsetmacro{\vertymin}{\verty - \zoomr}
\pgfmathsetmacro{\vertymax}{\verty + \zoomr}

\begin{tikzpicture}
  % No tick labels: the curve passes through all four of (+-1,0), (0,+-1).
  \begin{axis}[basic curve, width={180pt},
    xmin=-1.5, xmax=1.5, ymin=-1.5, ymax=1.5,
    xticklabels={}, yticklabels={},
    % Lift the title clear of the y-axis label, which sits at the top middle.
    title={The entire curve of \(x^{2} - xy + y^{2} = 1\)}, title style={yshift=1.5ex}]
    \addplot[very thick, domain=0:2*pi] ({cos(x) - sin(x)/sqrt(3)}, {cos(x) + sin(x)/sqrt(3)});
    \draw[main, thick] (axis cs:\zoomxmin,\zoomymin) rectangle (axis cs:\zoomxmax,\zoomymax);
    \fill[main] (axis cs:\zoomx,\zoomy) circle[radius=2pt] node[above left, main, font=\scriptsize] {\(P\)};
    \draw[supp, thick] (axis cs:\vertxmin,\vertymin) rectangle (axis cs:\vertxmax,\vertymax);
    \fill[supp] (axis cs:\vertx,\verty) circle[radius=2pt] node[left, supp, font=\scriptsize] {\(Q\)};
  \end{axis}
\end{tikzpicture}

% The zoom near (-1/sqrt(3), 1/sqrt(3)), titled #1, with the point labelled #2
% (may be empty); drawn as page 8 and again, retitled and labelled, as page 10.
\newcommand{\zoomregular}[2]{%
  \begin{tikzpicture}
    \begin{axis}[basic curve, width={180pt},
      axis lines=box, axis line style={main, thick},
      xlabel={\empty}, ylabel={\empty},
      xmin=\zoomxmin, xmax=\zoomxmax, ymin=\zoomymin, ymax=\zoomymax,
      xtick distance={0.2}, ytick distance={0.2},
      title style={align=center}, title={#1}]
      \addplot[very thick, domain={pi/2-0.5}:{pi/2+0.5}] ({cos(x) - sin(x)/sqrt(3)}, {cos(x) + sin(x)/sqrt(3)});
      \fill[main] (axis cs:\zoomx,\zoomy) circle[radius=2pt] node[above left, main] {#2};
    \end{axis}
  \end{tikzpicture}%
}

\zoomregular{Zoomed in near \(\left(-\tfrac{1}{\sqrt{3}}, \tfrac{1}{\sqrt{3}}\right)\)}{}

% Page 9: the same 5x zoom as page 8, but near Q.  However far we zoom in, the
% curve still doubles back there.  Framed in supp to match Q's box on page 7.
\begin{tikzpicture}
  \begin{axis}[basic curve, width={180pt},
    axis lines=box, axis line style={supp, thick},
    xlabel={\empty}, ylabel={\empty},
    xmin=\vertxmin, xmax=\vertxmax, ymin=\vertymin, ymax=\vertymax,
    xtick distance={0.2}, ytick distance={0.2},
    title style={align=center},
    title={\(y\) IS NOT an implicit function\\near \(Q = \left(-\tfrac{2}{\sqrt{3}}, -\tfrac{1}{\sqrt{3}}\right)\)}]
    \addplot[very thick, domain={5*pi/6-1}:{5*pi/6+1}] ({cos(x) - sin(x)/sqrt(3)}, {cos(x) + sin(x)/sqrt(3)});
    \fill[supp] (axis cs:\vertx,\verty) circle[radius=2pt] node[left, supp] {\(Q\)};
  \end{axis}
\end{tikzpicture}

% Page 10: page 8 retitled, to pair with page 9 after the definition.
\zoomregular{\(y\) IS an implicit function\\near \(P = \left(-\tfrac{1}{\sqrt{3}}, \tfrac{1}{\sqrt{3}}\right)\)}{\(P\)}
\end{document}
